% ============ 第4章 §4.5 Hilbert 变换的几乎处处收敛 (原书 pp.91-94) ============

\section{Hilbert 变换的几乎处处收敛}\label{sec:4.5}

$1 \leq p < \infty$, $f \in L^{p}(\mathbb{R})$ 时, Hilbert 变换 $H$ 由稠密性连续地扩张, 对这样扩张的 $H$, 考察 $Hf(x)$ 与
\[
\lim_{\varepsilon \downarrow 0}\int_{|y| > \varepsilon}\frac{f(x-y)}{y}\dd y
\]
的关系. $L^{2}$ 收敛的函数列取子列后也在逐点收敛的意义下收敛, 所以 $f \in C_{\mathrm{c}}^{\infty}(\mathbb{R})$ 的情形, 对几乎所有的 $x \in \mathbb{R}$ 这些一致. 这里证明 $f \in L^{p}(\mathbb{R})$, $1 \leq p < \infty$ 时, 对几乎所有的 $x \in \mathbb{R}$ 它们一致.

这里, 至今得到的成果被总动员.

(1) $1 < p < \infty$ 时, Hardy--Littlewood 极大算子 $M$ 与 Hilbert 变换 $H$ 在 $L^{p}(\mathbb{R})$ 中有界.

(2) $M$ 与 $H$ 在 $L^{1}(\mathbb{R})$ 中弱有界, 因此, 例如对一切 $\lambda > 0$ 与 $f \in L^{1}(\mathbb{R})$, $|\{Mf > \lambda\}| \lesssim \lambda^{-1}\|f\|_{L^{1}}$ 成立.

(3) 证明 $H$ 的弱 (1,1) 有界性时使用的 Calderón--Zygmund 分解这一技法.

(4) 证明 Lebesgue 微分定理时使用的稠密性讨论.

为了证明几乎处处极限存在, 考察
\[
H^{*}f(x) = \sup_{\varepsilon > 0}\left|\int_{|y| > \varepsilon}\frac{f(x-y)}{y}\dd y\right|
\]
具有什么样的性质是有效的. 作为关于 $H^{*}$ 的性质的第一条, 证明下面的不等式.

\begin{theorem}\label{thm:4.7}
对 $f \in C_{\mathrm{c}}^{\infty}(\mathbb{R})$, 对 $x \in \mathbb{R}$ 的逐点估计
\begin{equation}\tag{4.32}
H^{*}f(x) \leq M[Hf](x) + CMf(x)
\label{eq:4.32}
\end{equation}
成立, 因此, 对一切 $f \in L^{p}(\mathbb{R})$, $1 < p < \infty$,
\begin{equation}\tag{4.33}
\|H^{*}f\|_{L^{p}} \leq C\|f\|_{L^{p}}
\label{eq:4.33}
\end{equation}
成立.
\end{theorem}

若承认式~\eqref{eq:4.32}, 则式~\eqref{eq:4.33} 由 $M$, $H$ 都 $L^{p}$ 有界是显然的.

\begin{proof}
需要对每个 $\varepsilon > 0$ 证明
\[
\left|\int_{|y| > \varepsilon}\frac{f(x-y)}{y}\dd y\right| \leq M[Hf](x) + CMf(x).
\]
由于 $M[f(\varepsilon\cdot)] = Mf(\varepsilon\cdot)$, $H[f(\varepsilon\cdot)] = Hf(\varepsilon\cdot)$, 作变量替换只需考虑 $\varepsilon = 2$ 的情形.

令
\[
\tau(x) = \frac{15}{16}\chi_{(-1,1)}(x)(x^{2}-1)^{2} \quad (x \in \mathbb{R}).
\]
由积分公式
\[
\int_{\alpha}^{\beta}(x-\alpha)^{m}(\beta-x)^{n}\dd x = \frac{m!n!(\beta-\alpha)^{m+n+1}}{(m+n+1)!} \quad (\alpha, \beta \in \mathbb{R}, m, n \in \mathbb{N})
\]
可知 $\|\tau\|_{L^{1}} = 1$, 注意这一点. 于是,
\[
\int_{|y| > 2}\frac{f(x-y)}{y}\dd y = \int_{\mathbb{R}}\frac{f(x-y)}{y}\chi_{\mathbb{R}\setminus(-2,2)}(y)\dd y,
\]
所以
\begin{align*}
\int_{|y| > 2}\frac{f(x-y)}{y}\dd y - f * H\tau(x) &= \int_{\mathbb{R}}f(x-y)\left(\frac{\chi_{\mathbb{R}\setminus(-2,2)}(y)}{y} - H\tau(y)\right)\dd y.
\end{align*}
这里, 令
\[
\Psi(y) = \frac{\chi_{\mathbb{R}\setminus(-2,2)}(y)}{y} - H\tau(y), \quad \Phi(y) = \frac{1}{1+y^{2}}.
\]
$|x| > 3$ 时, 由支集的关系积分范围只需考虑 $[-1, 1]$, 所以
\[
|H\tau(x)| \leq \int_{\mathbb{R}}\frac{\tau(y)}{|y-x|}\dd y \leq \frac{1}{2}\|\tau\|_{L^{1}} = \frac{1}{2}.
\]
因此, 由于 $\tau$ 是具有紧支集的可微函数, 由式~\eqref{eq:4.7}, $\Psi \in L^{\infty}(\mathbb{R})$. 另外, $|y| > 2$ 时, $\chi_{\mathbb{R}\setminus(-2,2)}(y) = 1$, 所以
\begin{align*}
\Psi(y) &= \frac{1}{y} - \int_{\mathbb{R}}\frac{\tau(z)}{y-z}\dd z \\
&= \int_{\mathbb{R}}\left(\frac{\tau(z)}{y} - \frac{\tau(z)}{y-z}\right)\dd z \\
&= -\int_{\mathbb{R}}\frac{z\tau(z)}{y(y-z)}\dd z \\
&= O(y^{-2})
\end{align*}
成立. 因此, 存在常数 $C$ 使得 $|\Psi(y)| \leq \frac{C}{1+y^{2}} = C\Phi(y)$. 由这个不等式,
\begin{align*}
\left|\int_{|y| > \varepsilon}\frac{f(x-y)}{y}\dd y\right| &\leq |f * H\tau(x)| + \left|\int_{|y| > \varepsilon}\frac{f(x-y)}{y}\dd y - f * H\tau(x)\right| \\
&\leq |f * H\tau(x)| + C\Phi * |f|(x)
\end{align*}
可知成立. 由引理~\ref{lem:3.1} 与 $\|\Phi\|_{1} = \pi$,
\begin{align*}
|f * H\tau(x)| &= |Hf * \tau(x)| \leq M[Hf](x), \\
\Phi * |f|(x) &\leq \pi Mf(x)
\end{align*}
成立. 代入以上即得结论.
\end{proof}

这个命题是 $p > 1$ 时对 $L^{p}$ 函数的性质, $p = 1$ 时下面的结果成立.

\begin{theorem}\label{thm:4.8}
对一切 $\lambda > 0$ 与 $f \in L^{1}(\mathbb{R})$,
\begin{equation}\tag{4.34}
|\{H^{*}f > \lambda\}| \lesssim \frac{1}{\lambda}\|f\|_{L^{1}}
\label{eq:4.34}
\end{equation}
成立.
\end{theorem}

\begin{proof}
将 $f$ 按高度 $\lambda$ 作 Calderón--Zygmund 分解, 表为 $f = g + \sum_{j \in J}b_{j} = g + b$. 此时, 如练习题~\ref{exer:4.12} 所述, $H^{*}f \leq H^{*}g + H^{*}b$ 成立. 转到补集考虑显然地,
\begin{equation}\tag{4.35}
\{H^{*}f > \lambda\} \subset \left\{H^{*}g > \frac{\lambda}{2}\right\} \cup \left\{H^{*}b > \frac{\lambda}{2}\right\}
\label{eq:4.35}
\end{equation}
成立. 因此, 只需证明
\begin{equation}\tag{4.36}
\left|\left\{H^{*}g > \frac{\lambda}{2}\right\}\right| \lesssim \frac{1}{\lambda}\|f\|_{L^{1}},
\label{eq:4.36}
\end{equation}
\begin{equation}\tag{4.37}
\left|\left\{H^{*}b > \frac{\lambda}{2}\right\}\right| \lesssim \frac{1}{\lambda}\|f\|_{L^{1}}
\label{eq:4.37}
\end{equation}
这两个. 式~\eqref{eq:4.36} 容易证明. 与 $H$ 的弱 (1,1) 有界性的证明一样, 用 Chebyshev 不等式,
\[
|\{H^{*}g > \lambda\}| \leq \frac{1}{\lambda^{2}}\int_{\mathbb{R}}H^{*}g(x)^{2}\dd x
\]
开始. 依次使用定理~\ref{thm:4.7}, 式~\eqref{eq:4.17}, 式~\eqref{eq:4.16},
\[
|\{H^{*}g > \lambda\}| \leq \frac{1}{\lambda^{2}}\int_{\mathbb{R}}|g(x)|^{2}\dd x \lesssim \frac{1}{\lambda}\|g\|_{L^{1}} \lesssim \frac{1}{\lambda}\|f\|_{L^{1}}
\]
成立. 但确证式~\eqref{eq:4.37} 稍难. 实际上式~\eqref{eq:4.37} 如下证明. 首先, 与 $2I_{j}$ 同样地定义 $3I_{j}$, 由式~\eqref{eq:4.14},
\begin{equation}\tag{4.38}
\left|\bigcup_{j \in J}3I_{j}\right| \leq 3\sum_{j \in J}|I_{j}| = 3|\{M_{\dyadic}f > \lambda\}| \leq \frac{3}{\lambda}\|f\|_{L^{1}},
\label{eq:4.38}
\end{equation}
所以只需证明
\begin{equation}\tag{4.39}
\left|\{H^{*}b > \lambda\} \setminus \bigcup_{j \in J}3I_{j}\right| \lesssim \frac{1}{\lambda}\|f\|_{L^{1}}.
\label{eq:4.39}
\end{equation}
固定 $x \in \mathbb{R}\setminus\bigcup_{j \in J}3I_{j}$ 与 $\varepsilon > 0$.

(1) 对满足 $[x-\varepsilon, x+\varepsilon] \supset I_{j}$ 的 $j \in J$, $H_{\varepsilon}b_{j}(x) = 0$ 成立.

(2) 对满足 $[x-\varepsilon, x+\varepsilon] \cap I_{j} = \emptyset$ 的 $j \in J$, $H_{\varepsilon}b_{j}(x) = Hb_{j}(x)$ 成立.

(3) 设不满足这两个条件的 $j$ 的全体为 $Z$. 则对 $j \in Z$, $x - \varepsilon \notin I_{j}$ 且 $x + \varepsilon \in \overline{I_{j}}$, 或 $x - \varepsilon \in \overline{I_{j}}$ 且 $x + \varepsilon \notin I_{j}$ 成立. 因此, $\{I_{j}\}_{j=1}^{\infty}$ 不相交,
\[
I_{j} \subset (x-10\varepsilon, x+10\varepsilon) \quad (j \in Z),
\]
且对 $j \in Z$, $y \in I_{j}$, 若 $x - \varepsilon \in \overline{I_{j}}$, 则 $|c(I_{j}) - y| \lesssim \varepsilon$, $|x - y| \gtrsim \varepsilon$, $|x - c(I_{j})| \gtrsim \varepsilon$, 所以
\begin{align*}
\sum_{j \in Z}|H_{\varepsilon}b_{j}(x)| &\leq \frac{C}{\varepsilon}\sum_{j \in Z}\int_{I_{j}}|b_{j}(y)|\dd y \\
&\leq \frac{C}{\varepsilon}\int_{(x-10\varepsilon, x+10\varepsilon)}|f(y)|\dd y \\
&\leq CMf(x)
\end{align*}
成立.

因此, 式~\eqref{eq:4.39} 归结为
\[
\left|\left\{\sum_{j \in J}|Hb_{j}| + CMf > \lambda\right\} \setminus \bigcup_{j \in J}3I_{j}\right| \lesssim \frac{1}{\lambda}\|f\|_{L^{1}}.
\]
到了这一步, 与通常的 $H$ 的弱 (1,1) 有界性的证明相同.
\end{proof}

\begin{theorem}\label{thm:4.9}
$1 \leq p < \infty$, $f \in L^{p}(\mathbb{R})$ 时, 对几乎所有的 $x \in \mathbb{R}$, $H_{\varepsilon}f(x)$ 在 $\varepsilon \downarrow 0$ 时的极限存在.
\end{theorem}

\begin{proof}
$p > 1$ 的情形与 $p = 1$ 的情形类似, 所以证明 $p = 1$ 的情形. 只需证明极限不存在的点全体组成的集合 $E$ 的测度为 $0$. 于是, 用上极限 $\limsup$ 给 $E$ 的各点赋予级别, 将 $E$ 按级别分开考虑. 即, 对一切自然数 $k \in \mathbb{N}$,
\[
E_{k} \equiv \left\{x \in \mathbb{R} : \limsup_{\varepsilon, \varepsilon' \downarrow 0}|H_{\varepsilon}f(x) - H_{\varepsilon'}f(x)| > \frac{1}{k}\right\},
\]
则
\[
E = \bigcup_{k=1}^{\infty}E_{k},
\]
所以只需证明各 $E_{k}$ 的测度为 $0$. 注意, 极限的存在性对 $g \in C_{\mathrm{c}}^{\infty}(\mathbb{R})$ 是显然的. 任意取 $\varepsilon > 0$. 通过证明 $|E_{k}| \leq Ck\varepsilon$ 来完成定理的证明.

取 $g \in C_{\mathrm{c}}^{\infty}(\mathbb{R})$ 使得 $\|f - g\|_{L^{1}} \leq \varepsilon$. 这由稠密性是可能的. 作为集合的包含关系
\begin{equation}\tag{4.40}
E_{k} = \left\{x \in \mathbb{R} : \limsup_{\varepsilon, \varepsilon' \downarrow 0}|H_{\varepsilon}[f-g](x) - H_{\varepsilon'}[f-g](x)| > \frac{1}{k}\right\} \subset \left\{x \in \mathbb{R} : 2H^{*}[f-g](x) > \frac{1}{k}\right\}
\label{eq:4.40}
\end{equation}
成立, 所以由弱 (1,1) 不等式,
\[
|E_{k}| \leq Ck\|f - g\|_{L^{1}} \leq Ck\varepsilon.
\]
由于 $\varepsilon > 0$ 任意, 得 $|E_{k}| = 0$. 由于 $k \in \mathbb{N}$ 也任意, 得 $|E| = 0$. 即, 对几乎所有的 $x \in \mathbb{R}$, $H_{\varepsilon}f(x)$ 在 $\varepsilon \downarrow 0$ 时的极限存在.
\end{proof}

\begin{exercise}\label{exer:4.12}
对 $f, g \in L^{1}(\mathbb{R}^{n})$, 证明 $H^{*}[f+g](x) \leq H^{*}f(x) + H^{*}g(x)$ 成立.
\end{exercise}
